% =========================================================
% Applications in Quantum Technology
% =========================================================

\subsection{Micro Atomic Clocks}

The micro-scale trench and pillar structures fabricated in the MICRO experience serve as the mechanical and fluidic components of micro-fabricated atomic clocks. These are miniaturised versions of the caesium or rubidium vapour-cell atomic clocks that provide primary frequency standards. The deep reactive ion etched cavities form the sealed chambers that contain the alkali metal vapour and the optical interrogation components (a vertical-cavity surface-emitting laser (VCSEL) and a photodetector). Filter structures that function as \textit{atom blockers}, permitting gas molecules to pass whilst excluding larger particles, are also realised by deep silicon etching, exploiting the molecular-filter effect of nano-aperture arrays. As Italy's national metrology institute, INRIM directly researches these devices as the next generation of portable, chip-scale primary frequency standards.

\subsection{Plasmonic Nano-antennas}

The EBL-defined split-ring resonators produced in the NANO experience support \textit{localised surface plasmon resonances}: collective oscillations of conduction electrons at the metal surface driven by the incident electromagnetic field. The resonance frequency is determined by the geometry: the ring diameter, gap size, and metal thickness. By tuning these dimensions through EBL dose and etch depth, the antenna can be engineered to resonate at specific wavelengths in the visible or near-infrared spectrum.

The vertical nano-pillar geometry is of particular relevance to quantum technology applications. When the pillars are metallised and coupled to a quantum-well active medium (e.g.\ GaAs/AlGaAs), the antenna creates an optical nano-cavity around the emitter. The resulting enhancement of the local density of optical states accelerates spontaneous emission via the \textit{Purcell effect}, a direct quantum electrodynamical phenomenon exploited in single-photon sources, quantum repeaters, and quantum communication links.

\subsection{Nano-biosensors}

The same nano-antenna platform, when integrated on a microfluidic substrate, constitutes a sensitive nano-biosensor. The near-field enhancement at the split-ring gap concentrates the electromagnetic field by several orders of magnitude relative to the incident intensity, making the sensor highly responsive to refractive index changes caused by molecular binding events at the metal surface. Applications include label-free detection of proteins, nucleic acids, and small molecules relevant to drug-delivery monitoring and point-of-care diagnostics.
